\subsection{Plaintext model fidelity}

\begin{table*}[!t]\centering\small
\caption{Released DNAGPT evaluated locally on three genomic task families. Published comparators
are DNAGPT-M for signal recognition and abundance regression, and DNABERT-2 for the three
understanding-benchmark tasks. Signal recognition and abundance regression use released fine-tuned
heads; the understanding-benchmark rows use a locally fine-tuned head, as no head was released.}
\label{tab:baseline}
\begin{tabularx}{\textwidth}{Y l l l l}
\toprule
Task & Metric & This work & Reference & Source \\
\midrule
Polyadenylation-signal recognition$^{\ddagger}$ & accuracy & 0.9124 & 0.9151 & \cite{dnagpt2023} \\
 & F1 & 0.916 & --- & \\
Core promoter detection & MCC & 0.680 & 0.69 & \cite{dnabert2} \\
300\,bp promoter detection & MCC & 0.897 & 0.87 & \cite{dnabert2} \\
Splice-site classification & MCC & 0.831 & 0.85 & \cite{dnabert2} \\
mRNA abundance regression$^{\dagger}$ & $r^2$ & 0.562 & 0.62 & \cite{dnagpt2023} \\
 & Pearson $r$ & 0.753 & --- & \\
\bottomrule
\end{tabularx}
\vspace{2pt}
{\footnotesize $^{\ddagger}$The reference is DNAGPT's own reported accuracy for this model class,
measured on a held-out quarter of the set; this work evaluates all $22{,}604$ examples. The
comparison is same-model but not same-split.\\
$^{\dagger}$Model-fidelity evidence only. This task is outside the encrypted
scope; see Section~\ref{sec:limitations}.}
\end{table*}

Table~\ref{tab:baseline} and Fig.~\ref{fig:baseline} show that the target carries useful signal across classification and
regression tasks. On the complete polyadenylation-signal test set of $22{,}604$ examples, the
released classifier reaches $0.9124$ accuracy and $0.916$ F1, against the $0.9151$ accuracy
reported for this model class on a held-out quarter of that set~\cite{dnagpt2023}. The locally
fine-tuned understanding-benchmark heads reach MCC $0.680$, $0.897$, and $0.831$ on core-promoter,
promoter, and splice-site detection, against the $0.69$, $0.87$, and $0.85$ reported on the same
splits~\cite{dnabert2}. The released abundance-regression head reaches $r^2=0.562$ and Pearson
$r=0.753$, against the $0.62$ result reported by DNAGPT on the Xpresso dataset (Xpresso's own
published human $r^2$ is $0.59$~\cite{xpresso})~\cite{dnagpt2023}. These results establish model fidelity; encrypted evaluation is still
required to establish cryptographic feasibility.

Two dataset qualifications belong with these baseline results. DNAGPT does not release a head for
the understanding benchmark, so we trained a linear head and evaluated the three human promoter
and splice-site datasets selected for this study. For abundance regression, the historical source
was recovered and its split procedure reproduced, but the resulting test set may not be
gene-for-gene identical to the historical split. The abundance result is therefore model-fidelity
evidence and is not used as an encrypted target.

\begin{figure}[!t]\centering
  \includegraphics[width=\columnwidth]{fig_baseline.pdf}
  \caption{Plaintext model fidelity against published references. The signal-recognition comparison
  uses different test splits; promoter and splice-site comparisons use the stated benchmark splits.
  Abundance regression is outside the encrypted scope.}
  \label{fig:baseline}
\end{figure}
